\section{Reading the Lean proofs}\label{app:reading-lean-proofs}

Lean writes \(f\mathbin{\gg}g\) for the composite written \(gf\) in the paper.
In particular, \(\theta_X\mathbin{\gg}\sigma_S\) first forms an atomic sentence and then translates it.

\newcommand{\CEIproof}[3]{\href{https://github.com/flengyel/CategoricalEssentialIncompleteness/blob/a65aeb4d5aae48ef9d71bd3352f8935e7612e434/Lean/#1.lean\##2}{#3}}

\subsection{Restriction and extension of an interpretation}

In Theorem~\ref{thm:free-first-order-adjunction}, write $A=(X,a,b)$
for an object of $\mathsf{Pair}_c(\mathcal D)$ and $B=(P,T)$ for a consistent theory in $\mathsf{FO}_V(\mathcal D)$. The bijection
\[
 \operatorname{Hom}_{\mathsf{FO}_V(\mathcal D)}(\Sigma_{\mathrm{FO}} A,B)
 \cong
 \operatorname{Hom}_{\mathsf{Pair}_c(\mathcal D)}(A,J_{\mathrm{FO}}B)
\]
sends a translation $\sigma$ to the composite
$X\xrightarrow{\theta_X}S_X\xrightarrow{\sigma_S}S_P$.
Thus \texttt{restrict} is composition with the atomic sentence morphism. In the opposite direction,
\texttt{extend} takes a total morphism $f:X\to S_P$ satisfying
$a\leq f^*T$ and $b\leq f^*(\nu_P^*T)$. It interprets the generating
predicates by the sentences assigned by $f$. The variables retain their interpretation, and the assigned sentences are closed, so substitution preserves the assignment. The two inequalities put the images of the positive and negative axioms in $T$.
Deductive closure then gives a morphism of theories.

The two inverse identities in \texttt{homEquiv} use different facts.
Restricting $\sigma$ and then extending gives a translation that agrees
with $\sigma$ on each generating predicate. Uniqueness of interpretation
makes the two translations equal. Extending $f$ and then restricting
gives back $f$ by the interpretation rule for atomic sentences. 
Restricting the identity of $\Sigma_{\mathrm{FO}}A$ gives
the unit $\theta_X$. Extending the identity of $J_{\mathrm{FO}}B$
gives the counit, which interprets the predicate indexed by a sentence
$\varphi$ as $\varphi$ itself.\footnote{\CEIproof{CategoricalEssentialIncompleteness/V46SignedFirstOrderAdjunction}{L27-L56}{Restriction, extension and the inverse identities},
\CEIproof{CategoricalEssentialIncompleteness/V46SignedFirstOrderAdjunction}{L58-L91}{naturality, unit and counit}.}

\subsection{Separators and complete extensions}

Suppose $T_{a,b}$ is consistent, as in
Corollary~\ref{cor:signed-first-order-incompleteness}. In the proof from
inseparability, \texttt{intro U hU hc hcomp} introduces a closed
extension $U$, the inclusion $T_{a,b}\leq U$, its consistency and its
supposed completeness. Consistency makes $U$ and $\nu_X^*U$ disjoint.
Completeness makes their join the whole sentence object. Their inverse
images along $\theta_X$ are therefore complementary and contain $a$ and
$b$, respectively. This contradicts inseparability.

Conversely, a separator gives a morphism $d:X\to 1+1$ true on $a$ and
false on $b$. Interpret the predicates by $d$ on a one-element domain.
Equality is always true, and quantifying a formula leaves its value
unchanged. The true sentences form the theory \texttt{booleanTheory}.
If \texttt{h} asserts essential incompleteness, its second conjunct \texttt{h.2} says that every consistent closed extension of $T_{a,b}$ is incomplete. Applied to \texttt{booleanTheory}, it contradicts completeness.\footnote{\CEIproof{CategoricalEssentialIncompleteness/V46SignedSubobjectIncompleteness}{L43-L75}{The two implications},
\CEIproof{CategoricalEssentialIncompleteness/V46SignedBooleanExtension}{L28-L74}{the theory of the Boolean valuation}. The Palomar proof source gives the \CEIproof{CEIPalomar/Solution}{L83-L110}{equivalence for subobjects and the Boolean assignment}.}

\subsection{One valuation for two derivations}

Lemma~\ref{lem:arithmetic-signed-theory} supplies the consistency
hypothesis for the arithmetic diagonal pair. Suppose that derivations
$p$ and $q$ conclude $\varphi$ and $\neg\varphi$. The definition
\texttt{jointSupport} concatenates their lists of positive axiom labels.
Give $P_n$ the value true exactly when $n$ occurs in this list.
Every positive leaf of either derivation is true. Every negative leaf
has label in the negative diagonal domain, which is disjoint from the
positive domain, so the corresponding negated atom is true as well.

Induction on each coded derivation gives truth of its conclusion for the
same list. The line \texttt{rw [h] at hq} replaces the conclusion of
$q$ by the negation of the conclusion of $p$ in the second truth
assertion. The Boolean rule for negation now gives a contradiction.
The list depends on $p$ and $q$, so it need only account for their
leaves. A separator would have to assign the prescribed values at every
label in either diagonal domain. In the formal proof the list
and derivations are coded in an arbitrary model of $\PA$. Internal induction includes proof codes of nonstandard length when the model is nonstandard.
First-order completeness therefore gives a proof in $\PA$ that no such pair of
derivation codes exists.\footnote{\CEIproof{CategoricalEssentialIncompleteness/V46ArithmeticSignedContradiction}{L32-L56}{The common list and its two truth assertions},
\CEIproof{CategoricalEssentialIncompleteness/V46ArithmeticSignedContradiction}{L85-L92}{the contradiction},
\CEIproof{CategoricalEssentialIncompleteness/V46ArithmeticSignedConsistencyPA}{L45-L64}{the proof in PA}. The Palomar proof of the \CEIproof{CEIPalomar/Solution}{L114-L145}{existential consequence} takes the two diagonal subobjects in the predicate category of $\PA$ as its witness.}

\subsection{The lift of a conditional}

In Theorem~\ref{thm:comparison}, the conditional $s$ belongs to the
category in which programs are interpreted. The proof needs a morphism
$\ell:U\rightharpoonup U+U$ in the recursion category to which it can
apply the Turing property. Lean's \texttt{turing\_diagonal} takes
$\ell$ as an argument and gives a total $g:U\to U$ satisfying
$hg=\ell g$, where $h=b\tau\Delta_U$. By hypothesis $s$ is the interpretation of $\ell$, up to the coproduct isomorphism, and its two summand domains are $N,P$, in that order. Pulling back
the two given inclusions along $Fg$ therefore puts each of
$(Fg)^*P$ and $(Fg)^*N$ below the other. Antisymmetry, written \texttt{le\_antisymm},
gives $(Fg)^*P=(Fg)^*N$. Their meet is zero by disjointness, so each
is zero.\footnote{\CEIproof{CategoricalEssentialIncompleteness/V46Comparison}{L35-L89}{The comparison proof}.}
Proposition~\ref{prop:missing-lift} gives a conditional in partial
functions whose output summand decides membership in a nonrecursive
set. Every total source morphism in that example is recursive. Hence
the conditional has no lift.\footnote{\CEIproof{CategoricalEssentialIncompleteness/OracleMissingLift}{L89-L103}{The conditional has no source lift}.}

\subsection{Evaluation of a theory in the slice}

In Theorem~\ref{thm:associated-incompleteness}, let $u:U\hookrightarrow S$
be the theorem subobject. An extension $W$ of the constant family
$B_U(u)$ can vary over $U$. The fixed logical theorem determines a
point $\lambda_U:1\to U$. Pullback along this point evaluates $W$ at
one fiber. Lean's \texttt{evaluationTheory} pairs the resulting
sentence subobject with a proof of deductive closure. Evaluation transports the rule algebra, and initiality puts the conclusions of derivations in the evaluated subobject. If $W$ were consistent and complete, its
fiber would be consistent and complete. It would extend $u$ because
evaluation takes $B_U(u)$ back to $u$. Essential incompleteness of $u$ excludes such a fiber,
and hence excludes every consistent complete closed extension of
$B_U(u)$.\footnote{\CEIproof{CategoricalEssentialIncompleteness/V46AssociatedClosedTheories}{L68-L104}{Evaluation, closure, and the arbitrary slice extension}.}

\Needspace{9\baselineskip}
\subsection{One morphism for all consistent presentations}

In Theorem~\ref{thm:uniform-independent-sentences}, the order of the
first two proof steps gives the uniformity. Lean first supplies
\texttt{sentence} as the morphism $q:E\to S$. Only then does
\texttt{intro e hc} introduce the index $e$ and the hypothesis that
$\Gamma_e$ is consistent. For this $e$, the proof identifies the
numerical value $(G_{\N}q)(e)$ with the code of a closed formula.
The lemma \texttt{output\_avoids} uses consistency at $e$ to show that the diagonal output belongs to neither the positive nor the negative proof-search domain. A derivation of the formula would put the output in the positive domain. A derivation of its negation would put it in the negative domain.\footnote{\CEIproof{CategoricalEssentialIncompleteness/V46ArithmeticUniformIndependent}{L184-L218}{The morphism and the two excluded derivations}. See also the \CEIproof{CategoricalEssentialIncompleteness/V46ArithmeticUniformIndependent}{L116-L121}{avoidance of both proof-search domains}.}
